\documentclass[10pt,a4paper]{amsart}
\usepackage[margin=19mm]{geometry}
\usepackage{amsmath,amssymb,mathtools}
\usepackage[scr=rsfs,cal=euler]{mathalfa}
\usepackage{xfrac}
\usepackage[unicode,hidelinks,bookmarks=false]{hyperref}
\newcommand{\ie}{i.e.\ }
\newcommand{\cf}{cf.\ }
\newcommand{\resp}{respectively\ }
\newcommand{\Subsection}{\subsection}
\providecommand{\nobreakdash}{\nobreak\hbox{-}}
\setlength{\emergencystretch}{2em}
\multlinegap0pt
\usepackage{mathtools}
\usepackage{mathrsfs}
\usepackage{enumitem}
\usepackage{bm}
\usepackage{stmaryrd}
\usepackage{tikz-cd}
\usepackage{dsfont}
\hypersetup{
  colorlinks=true,
  linkcolor=blue,
  citecolor=blue,
  urlcolor=blue
}
\newtheorem{theorem}{Theorem}[section]
\newtheorem{proposition}[theorem]{Proposition}
\newtheorem{lemma}[theorem]{Lemma}
\newtheorem{corollary}[theorem]{Corollary}
\newtheorem{definition}[theorem]{Definition}
\newtheorem{remark}[theorem]{Remark}
\newtheorem{example}[theorem]{Example}
\newtheorem{assumption}[theorem]{Assumption}
\newcommand{\R}{\mathbb{R}}
\newcommand{\C}{\mathbb{C}}
\newcommand{\Rtimes}{\mathbb{R}^{\times}}
\newcommand{\dd}{\mathrm{d}}
\newcommand{\ii}{\mathrm{i}}
\newcommand{\id}{\mathrm{id}}
\newcommand{\Lie}{\mathcal{L}}
\newcommand{\ad}{\mathrm{ad}}
\newcommand{\Ad}{\mathrm{Ad}}
\newcommand{\Tr}{\mathrm{Tr}}
\newcommand{\End}{\mathrm{End}}
\newcommand{\Hom}{\mathrm{Hom}}
\newcommand{\Span}{\mathrm{span}}
\newcommand{\im}{\mathrm{im}}
\newcommand{\rank}{\mathrm{rank}}
\newcommand{\Ker}{\mathrm{Ker}}
\newcommand{\pr}{\mathrm{pr}}
\newcommand{\supp}{\mathrm{supp}}
\newcommand{\diag}{\mathrm{diag}}
\newcommand{\su}{\mathfrak{su}}
\newcommand{\g}{\mathfrak{g}}
\newcommand{\e}{\mathrm{e}}
\newcommand{\pair}[2]{\left\langle #1,#2\right\rangle}
\newcommand{\abs}[1]{\left|#1\right|}
\newcommand{\norm}[1]{\left\|#1\right\|}
\newcommand{\set}[1]{\left\{#1\right\}}
\newcommand{\br}[1]{\left(#1\right)}
\newcommand{\sq}[1]{\left[#1\right]}
\newcommand{\mc}[1]{\mathcal{#1}}
\newcommand{\Jone}{J^1}
\newcommand{\Lbundle}{\mathcal{L}}
\newcommand{\Pbundle}{P}
\newcommand{\Cphase}{\mathcal{C}}
\newcommand{\Pjet}{\mathcal{P}}
\newcommand{\Tkin}{\mathcal{T}}
\newcommand{\epsc}{\varepsilon^{\mathrm c}}
\newcommand{\FL}{\mathbb{F}L}
\begin{document}
\title[Contact Tulczyjew Geometry for Continuous and Discrete Dissi]{Contact Tulczyjew Geometry for Continuous and Discrete Dissipative Dynamics on Skew Algebroids}
\author{Leonardo Colombo, Manuel de León}
\date{}
\maketitle
\noindent\emph{Source and licence.} Contact Tulczyjew Geometry for Continuous and Discrete Dissipative Dynamics on Skew Algebroids. Version arXiv:2606.10515v1 (CC BY 4.0). This material is distributed under the Creative Commons Attribution 4.0 International licence, http://creativecommons.org/licenses/by/4.0/, as recorded in the arXiv OAI licence metadata for this exact version and shown on the abstract page https://arxiv.org/abs/2606.10515v1. Full rights notice: licenses/arxiv-2606.10515-CC-BY-4.0.txt.\par\smallskip
\noindent\textit{Editorial scope.} Counted unit: the original \texttt{theorem} environment \texttt{thm:implicitdynamics} at source lines 534--580 together with its complete proof at lines 582--650; the delivered document keeps source lines 332--945 so that every referenced label and every citation resolves inside it. Native editable source is the paper's own concatenated TeX; the AMS wrapper is editorial formatting only. No publisher class, upstream script or unrelated artwork is redistributed.
\begin{definition}\label{var}
Let \(a(t)\) be an admissible curve with base curve \(x(t)\). A
\emph{variation generator} along \(a(t)\) is a time-dependent section $\eta(t)=\eta^\alpha(t)e_\alpha$ along \(x(t)\). The associated infinitesimal admissible variation is given
locally by
\begin{equation}
        \delta x^i
        =
        \rho^i_{\alpha}(x)\eta^\alpha,
        \qquad
        \delta y^\gamma
        =
        \dot\eta^\gamma
        +
        C^\gamma_{\alpha\beta}(x)y^\alpha\eta^\beta .
\label{eq:admvar}
\end{equation}
\end{definition}

\begin{remark}
Formula \eqref{eq:admvar} is the standard local expression of admissible variations
on a general algebroid, following \cite{GG2008}.
It can also be described intrinsically through the relation $\kappa_E$
associated with the Tulczyjew morphism $\varepsilon_E$, but only the
local form \eqref{eq:admvar} will be used below. Thus the admissible
variations are not introduced as an additional choice; they are determined by the algebroid structure itself. For variational problems with fixed
endpoints we shall impose $\eta(t_0)=\eta(t_1)=0$.
\end{remark}

\subsection{Contact Tulczyjew structures and local trivializations}
\label{sec:contact-structure-local}

The classical Tulczyjew picture is formulated in terms of canonical symplectic
structures on iterated tangent and cotangent bundles \cite{Tulczyjew1976}. In the
contact case, the appropriate geometric language is not that of a chosen global
contact form, but that of line bundles and homogeneous
\(\mathbb R^\times\)-principal bundles. Following the contact Tulczyjew
construction of \cite{GG2024}, the contact model is expressed in terms of first
jets of sections of a line bundle, rather than through a globally trivial
contactization \(T^*M\times\mathbb R\).

Let \(E \to M\) be an skew algebroid. Following the line-bundle approach to
contact geometry \cite{GG2024}, the contact structure on the phase side is encoded
by a line bundle $\mathcal L \rightarrow E^*$ or, equivalently, by the associated principal \(\mathbb R^\times\)-bundle $P=\mathcal L^\times\rightarrow E^*$.
After choosing a local trivialization of \(\mathcal L\), the line-bundle coordinate
is represented by an ordinary scalar variable \(z\). Thus the local contact phase
space is represented by \(E^*\times\mathbb R\), with coordinates
\((x^i,p_\alpha,z)\). The coordinate \(z\) is not intrinsic; it depends on the
chosen trivialization.

Similarly, a contact generating object is intrinsically a line-bundle
object. In a local trivialization it is represented by an ordinary smooth
function
\[
        L:E\times\R\rightarrow \R,
        \qquad
        (x,y,z)\mapsto L(x,y,z).
\]


In the same local trivialization, the source of the contact Tulczyjew morphism is represented by coordinates $(x^i,y^\alpha,z,p_i,p_\alpha,p_z,u)$, where \((p_i,p_\alpha)\) are the cotangent variables dual to
\((x^i,y^\alpha)\), \(p_z\) is the contact momentum, and \(u\) is the
contact value variable. The target is represented by coordinates
$(x^i,p_\alpha,z,\dot x^i,\dot p_\alpha,\dot z)$, namely by contactified kinematic variables over \(E^*\times\R\).

The local contact Tulczyjew morphism is
\begin{equation}
\epsc(x,y,z,p_x,p_y,p_z,u)
=
\bigl(
 x,p_y,z,
 \rho(x)y,
 \rho(x)p_x-C(x)(y,p_y)+p_zp_y,
 u
\bigr).
\label{eq:epscmodel}
\end{equation}
Equivalently,
\[
        \dot x^i
        =
        \rho^i_{\alpha}(x)y^\alpha,
        \qquad
        \dot p_\alpha
        =
        \rho^i_{\alpha}(x)p_i
        -
        C^\gamma_{\alpha\beta}(x)y^\beta p_\gamma
        +
        p_zp_\alpha,
        \qquad
        \dot z=u .
\]
\begin{remark}
The term \(p_zp_\alpha\) is the local contact contribution. In
Section~\ref{sec:intrinsic} we show that, in the trivial
contactization, \eqref{eq:epscmodel} is obtained from the ordinary
algebroid Tulczyjew morphism \(\varepsilon_E:T^*E\to TE^*\) by adding the
Euler vector field contribution on \(E^*\).
\end{remark}


The local contact differential of \(L\) is represented by
\[
J^{1,c}L(x,y,z)
=
\left(
 x,y,z,
 \frac{\partial L}{\partial x^i},
 \frac{\partial L}{\partial y^\alpha},
 \frac{\partial L}{\partial z},
 L
\right).
\]
The local contact Legendre map is
\begin{equation}
    \label{CLM}
        \lambda_L(x,y,z)
        =
        \left(
        x,\frac{\partial L}{\partial y},z
        \right).\end{equation}

Applying \(\epsc\) to \(J^{1,c}L\), we obtain the generated phase relation
\[
\Lambda_L(x,y,z)
=
\left(
 x,\frac{\partial L}{\partial y},z,
 \rho(x)y,
 \rho(x)\frac{\partial L}{\partial x}
 -
 C(x)\left(y,\frac{\partial L}{\partial y}\right)
 +
 \frac{\partial L}{\partial z}\frac{\partial L}{\partial y},
 L
\right).
\]

Let $\gamma(t)=(x(t),y(t),z(t))$ be a curve in \(E\times\R\). The local contact Legendre map sends
\(\gamma\) to the phase curve
\[
        \lambda_L\circ\gamma
        =
        \left(
        x(t),
        \frac{\partial L}{\partial y}(x(t),y(t),z(t)),
        z(t)
        \right)
        \in E^*\times\R .
\]
We shall call $\lambda_L\circ\gamma$ the local Legendre curve associated with \(\gamma\).
We denote by $\mathbf t(\lambda_L\circ\gamma)$ its tangent
prolongation, written in local coordinates as
\begin{equation}\label{tangprolong}
 \mathbf t(\lambda_L\circ\gamma)
        =
        \left(
        x,
        \frac{\partial L}{\partial y},
        z,
        \dot x,
        \frac{d}{dt}\frac{\partial L}{\partial y},
        \dot z
        \right).
\end{equation}
The contact Tulczyjew dynamics is defined by the matching condition
\begin{equation}
        \mathbf t(\lambda_L\circ\gamma)
        =
        \Lambda_L\circ\gamma .
\label{eq:kinematicmatch}
\end{equation}
Equating the components in \eqref{eq:kinematicmatch} gives the local
implicit equations derived in the next section.

\begin{remark}
All formulas in this subsection are local representatives of the
\(\mathcal L\)-line-bundle picture. The nontrivial-line-bundle construction
requires the full homogeneous \(\mathbb R^\times\)-principal-bundle formulation. In
Section~\ref{sec:intrinsic} we prove the canonical origin of the above local morphism in the
trivial contactization.
\end{remark}

\section{Implicit dissipative dynamics}
\label{sec:dynamics}

We now derive the local dynamics determined by the contact Tulczyjew
matching condition \eqref{eq:kinematicmatch}.
The guiding point is the same as in the classical Tulczyjew formalism:
the dynamics is first a geometric relation on the phase space, and only
under additional regularity assumptions does it become the graph of a
vector field. Thus singular Lagrangians are not excluded from the
construction.

Let $L:E\times\R\longrightarrow\R$ be the local representative of a contact generating object in a chosen
trivialization of the line bundle. 

\begin{definition}
A curve $\gamma(t)=(x(t),y(t),z(t))\in E\times\R$ is called a \emph{contact Tulczyjew trajectory} for \(L\) if its
Legendre image \(\lambda_L\circ\gamma\) satisfies the matching condition $\mathbf t(\lambda_L\circ\gamma)=\Lambda_L\circ\gamma$.
\end{definition}


\begin{theorem}
\label{thm:implicitdynamics} $\gamma(t)=(x^i(t),y^\alpha(t),z(t))$ is a contact Tulczyjew
trajectory for \(L\) if and only if it satisfies
\begin{align}
\dot x^i
&=
\rho^i_{\alpha}(x)y^\alpha,
\label{eq:ELH1}
\\
\frac{\dd}{\dd t}
\left(
\frac{\partial L}{\partial y^\alpha}
\right)
&=
\rho^i_{\alpha}(x)\frac{\partial L}{\partial x^i}
-
C^\gamma_{\alpha\beta}(x)y^\beta
\frac{\partial L}{\partial y^\gamma}
+
\frac{\partial L}{\partial z}
\frac{\partial L}{\partial y^\alpha},
\label{eq:ELH2}
\\
\dot z
&=
L(x,y,z).
\label{eq:ELH3}
\end{align}
Equivalently,
\begin{equation}
\frac{\dd}{\dd t}
\left(
\frac{\partial L}{\partial y^\alpha}
\right)
+
C^\gamma_{\alpha\beta}(x)y^\beta
\frac{\partial L}{\partial y^\gamma}
-
\rho^i_{\alpha}(x)\frac{\partial L}{\partial x^i}
-
\frac{\partial L}{\partial z}
\frac{\partial L}{\partial y^\alpha}
=
0.
\label{eq:ELHfinal}
\end{equation}
\end{theorem}

\begin{proof}
Let \(\gamma(t)=(x(t),y(t),z(t))\). The Legendre image of \(\gamma\) is
\[
        \lambda_L\circ\gamma
        =
        \left(
        x(t),
        \frac{\partial L}{\partial y}(x(t),y(t),z(t)),
        z(t)
        \right).
\]
Hence its tangent prolongation is locally
\[
        \mathbf t(\lambda_L\circ\gamma)
        =
        \left(
        x,
        \frac{\partial L}{\partial y},
        z,
        \dot x,
        \frac{\dd}{\dd t}\frac{\partial L}{\partial y},
        \dot z
        \right).
\]
On the other hand,
\[
\Lambda_L\circ\gamma
=
\left(
 x,\frac{\partial L}{\partial y},z,
 \rho(x)y,
 \rho(x)\frac{\partial L}{\partial x}
 -
 C(x)\left(y,\frac{\partial L}{\partial y}\right)
 +
 \frac{\partial L}{\partial z}
 \frac{\partial L}{\partial y},
 L
\right)
\]
along the curve. Therefore the matching condition
$\mathbf t(\lambda_L\circ\gamma)
        =
        \Lambda_L\circ\gamma$ is equivalent, component by component, to
\[
        \dot x=\rho(x)y,
\]
\[
        \frac{\dd}{\dd t}
        \frac{\partial L}{\partial y}
        =
        \rho(x)\frac{\partial L}{\partial x}
        -
        C(x)\left(y,\frac{\partial L}{\partial y}\right)
        +
        \frac{\partial L}{\partial z}
        \frac{\partial L}{\partial y},
\]
and
\[
        \dot z=L(x,y,z).
\]
In coordinates these are precisely \eqref{eq:ELH1}-\eqref{eq:ELH3}.
Equivalently, writing the momentum equation with all terms on the left-hand side
gives \eqref{eq:ELHfinal}. Conversely, if these three component equations hold, then the tangent
prolongation of the Legendre curve and the generated phase relation have
the same coordinates along \(\gamma\). Hence
\(t(\lambda_L\circ\gamma)=\Lambda_L\circ\gamma\).
\end{proof}



\begin{definition}
A local representative \(L\) is called \emph{regular} if the fiber Hessian $\left(\frac{\partial^2L}{\partial y^\alpha\partial y^\beta}\right)$ is nondegenerate, and \emph{singular} otherwise.
\end{definition}




Regularity implies local invertibility of the contact Legendre map $\lambda_L:E\times\mathbb R\rightarrow E^*\times\mathbb R$. Indeed, since \(\lambda_L\) is the identity on the variables \((x,z)\), its local
invertibility is equivalent to the nondegeneracy of the fiber Hessian with respect
to \(y\).

If \(L\) is regular, equation (7) can be solved locally for \(\dot y\), because
\(\dot x\) and \(\dot z\) are already determined by (6) and (8). Hence the
implicit relation defined by (6)-(8) determines locally a vector field on
\(E\times\mathbb R\). If \(L\) is singular, equations (6)-(8) define, in general, an implicit
differential relation on \(E\times\mathbb R\), rather than the graph of a vector
field. Thus singularity is not an anomaly of the theory, but part of its natural
domain of applicability.

\subsection{Variational interpretation and Euler-Lagrange-Herglotz equations}
\label{sec:lagrangian}

We now relate the contact-Tulczyjew dynamics to the Herglotz variational
principle on the skew algebroid $E$. The admissibility condition and the
admissible variations are those introduced in Definition~2.3 and
Definition~2.5.

\begin{definition}
Let \(L:E\times\mathbb R\to\mathbb R\) be the local representative of a
contact generating object. A curve $\gamma(t)=(x(t),y(t),z(t))$
is called a Herglotz extremal if \((x(t),y(t))\) is admissible,
$\dot z=L(x,y,z)$, and the terminal value \(z(T)\) is stationary under all admissible
variations with fixed endpoints.
\end{definition}

\begin{theorem}
A curve \(\gamma(t)=(x(t),y(t),z(t))\) is a Herglotz extremal, with respect to
admissible variations with fixed endpoints and fixed initial value of \(z\), if
and only if it is a contact Tulczyjew trajectory.
\end{theorem}

\begin{proof}
Let \(\gamma(t)=(x(t),y(t),z(t))\) be a Herglotz extremal. Since
\(\dot z=L(x,y,z)\), the induced variation of \(z\) satisfies
\[
\frac{d}{dt}(\delta z)
=
\frac{\partial L}{\partial x^i}\delta x^i
+
\frac{\partial L}{\partial y^\alpha}\delta y^\alpha
+
\frac{\partial L}{\partial z}\delta z .
\]
Let
\[
\lambda(t)=
\exp\left(
-\int_0^t
\frac{\partial L}{\partial z}(\tau)\,d\tau
\right).
\]
Since the initial value of \(z\) is fixed, \(\delta z(0)=0\). Multiplying by
\(\lambda(t)\) gives
\[
\frac{d}{dt}\bigl(\lambda(t)\delta z(t)\bigr)
=
\lambda(t)
\left(
\frac{\partial L}{\partial x^i}\delta x^i
+
\frac{\partial L}{\partial y^\alpha}\delta y^\alpha
\right).
\]
Therefore, because \(\lambda(T)\neq 0\), the stationarity condition
\(\delta z(T)=0\) is equivalent to
\[
\int_0^T
\lambda(t)
\left(
\frac{\partial L}{\partial x^i}\delta x^i
+
\frac{\partial L}{\partial y^\alpha}\delta y^\alpha
\right)\,dt=0.
\]

Using the admissible variations of Definition \ref{var}, and integrating by parts, with \(\eta(0)=\eta(T)=0\), we obtain
\[
\int_0^T
\lambda(t)
\left[
\rho^i_\alpha\frac{\partial L}{\partial x^i}
-
C^\gamma_{\alpha\beta}y^\beta
\frac{\partial L}{\partial y^\gamma}
-
\frac{d}{dt}
\left(
\frac{\partial L}{\partial y^\alpha}
\right)
+
\frac{\partial L}{\partial z}
\frac{\partial L}{\partial y^\alpha}
\right]
\eta^\alpha\,dt=0 .
\]
Since the generators \(\eta^\alpha\) are arbitrary, the bracketed term vanishes.
Equivalently,
\[
\frac{d}{dt}
\left(
\frac{\partial L}{\partial y^\alpha}
\right)
+
C^\gamma_{\alpha\beta}y^\beta
\frac{\partial L}{\partial y^\gamma}
-
\rho^i_\alpha
\frac{\partial L}{\partial x^i}
-
\frac{\partial L}{\partial z}
\frac{\partial L}{\partial y^\alpha}
=0 .
\]
Together with the admissibility equation and the Herglotz equation for \(z\), this
is precisely the system characterized in Theorem~3.2. Hence \(\gamma\) is a
contact Tulczyjew trajectory.

Conversely, assume that \(\gamma\) is a contact Tulczyjew trajectory. Then, by
Theorem~3.2, it satisfies the admissibility equation, the Herglotz equation, and
the Euler--Lagrange--Herglotz equations above. Substituting these equations in the
preceding integration-by-parts identity gives
\[
\int_0^T
\lambda(t)
\left(
\frac{\partial L}{\partial x^i}\delta x^i
+
\frac{\partial L}{\partial y^\alpha}\delta y^\alpha
\right)\,dt=0
\]
for every admissible variation with fixed endpoints. Since
\[
\lambda(T)\delta z(T)
=
\int_0^T
\lambda(t)
\left(
\frac{\partial L}{\partial x^i}\delta x^i
+
\frac{\partial L}{\partial y^\alpha}\delta y^\alpha
\right)\,dt
\]
and \(\lambda(T)\neq 0\), it follows that \(\delta z(T)=0\). Thus \(\gamma\) is a
Herglotz extremal.
\end{proof}

\begin{remark}
Thus the local contact-Tulczyjew equations are not merely the coordinate
expression of the matching condition
$t(\lambda_L\circ\gamma)=\Lambda_L\circ\gamma$. They also encode the Herglotz variational principle on the skew algebroid.
In this sense, the Euler-Lagrange-Herglotz equations are the variational
interpretation of the contact Tulczyjew trajectory condition.
\end{remark}



\begin{definition}
The local representative $L$ is called hyperregular if the contact Legendre map
$\lambda_L:E\times\R\to E^*\times\R$ defined in \eqref{CLM} is a global diffeomorphism.
\end{definition}

\begin{proposition}
\label{prop:hyperregular}
Assume that \(L\) is hyperregular. Then there is a unique smooth map $Y:E^*\times\R\to E$, with \(Y(x,p,z)\in E_x\), determined by
$\lambda_L(x,Y(x,p,z),z)=(x,p,z)$, and hence a unique Hamiltonian \(H:E^*\times\R\to\R\) defined by
\[
H(x,p,z)=\langle p,Y(x,p,z)\rangle-L(x,Y(x,p,z),z).
\]
Under the contact Legendre map, the contact-Tulczyjew dynamics generated
by \(L\) is equivalent to the contact-Hamiltonian dynamics generated by
\(H\).
\end{proposition}

\begin{proof}
Since \(L\) is hyperregular, \(\lambda_L\) is a global diffeomorphism. Therefore,
for every \((x,p,z)\in E^*\times\R\), there is a unique
\(y=Y(x,p,z)\in E_x\) such that $p_\alpha=\frac{\partial L}{\partial y^\alpha}(x,y,z)$.
This proves that \(H\) is well defined and smooth.

Differentiating \(H\) with respect to \(p_\alpha\), \(x^i\), and \(z\), and using
\(p_\alpha=\partial L/\partial y^\alpha\), all terms containing derivatives of
\(Y\) cancel. Hence
\[
\frac{\partial H}{\partial p_\alpha}=Y^\alpha,\qquad
\frac{\partial H}{\partial x^i}
=
-\frac{\partial L}{\partial x^i}(x,Y,z),
\qquad
\frac{\partial H}{\partial z}
=
-\frac{\partial L}{\partial z}(x,Y,z).
\]

Let \(\gamma(t)=(x(t),y(t),z(t))\) be a contact Tulczyjew trajectory and set
$p_\alpha(t)=\frac{\partial L}{\partial y^\alpha}(x(t),y(t),z(t))$. By Theorem~\ref{thm:implicitdynamics},
\[
\dot x^i=\rho^i_\alpha(x)y^\alpha,
\]
\[
\dot p_\alpha
=
\rho^i_\alpha(x)\frac{\partial L}{\partial x^i}
-
C^\gamma_{\alpha\beta}(x)y^\beta p_\gamma
+
\frac{\partial L}{\partial z}p_\alpha,
\]
and $\dot z=L(x,y,z)$.

Using the identities above and \(y=Y(x,p,z)\), these equations become
\[
\dot x^i
=
\rho^i_\alpha(x)\frac{\partial H}{\partial p_\alpha},
\]
\[
\dot p_\alpha
=
-\rho^i_\alpha(x)\frac{\partial H}{\partial x^i}
-
C^\gamma_{\alpha\beta}(x)
\frac{\partial H}{\partial p_\beta}p_\gamma
-
\frac{\partial H}{\partial z}p_\alpha,
\]
and
\[
\dot z
=
p_\alpha\frac{\partial H}{\partial p_\alpha}-H.
\]
Thus the Lagrangian and Hamiltonian contact equations are equivalent under
the hyperregular Legendre transformation. Conversely, any solution of the
Hamiltonian equations pulls back by \(\lambda_L^{-1}\) to a solution of the
contact-Tulczyjew equations.
\end{proof}

\begin{remark}
 The singular case should be understood in the relational sense. The generated
object is the subset $\mathcal R_L:=\operatorname{Im}(\Lambda_L)\subset T(E^*\times \mathbb R)$, and the dynamics is defined by the matching condition $t(\lambda_L\circ\gamma)\in \mathcal R_L$. If the contact Legendre map is regular, this condition may be solved locally as
a vector field. If it is hyperregular, this vector field is transformed by the
Legendre map into the contact Hamiltonian dynamics described above. If the
Legendre map is singular, however, no such resolution is available in general,
and the system should be treated as an implicit differential relation \cite{deLeonLainz2019}.

The passage from an implicit differential relation to an actual space of integrable trajectories is a separate geometric problem. In the general theory of implicit differential equations, this is handled by an integrability or constraint algorithm: one projects the relation to the phase space, imposes the tangency condition, projects again, and iterates until an integrable final relation is reached; see the foundational work of Mendella, Marmo and Tulczyjew~\cite{MendellaMarmoTulczyjew1995}. We do not perform here the corresponding stabilization procedure in the contact algebroid setting. Extending such an algorithm to contact Tulczyjew dynamics on skew algebroids would require a precontact constraint theory compatible with algebroid admissibility and with the line-bundle contact structure. This is a separate problem, and is left for future work.
\end{remark}

\begin{remark}
Recall that, for a Lie algebroid \(E\to M\) and a fibration \(P\to M\), the
standard prolongation is
$\mathcal T^E P:=E\times_{TM}TP$,
with anchor given by the projection to \(TP\); see
\cite{Martinez2001}. Thus, in the regular Lie algebroid case, the relation
$\mathcal R_L\subset T(E^*\times\mathbb R)$
can be equivalently described as the image, under the anchor $\rho_\pi:\mathcal T^E(E^*\times\mathbb R)\rightarrow T(E^*\times\mathbb R)$, of the Hamiltonian section of the contact Hamiltonian prolongation
\(\mathcal T^E(E^*\times\mathbb R)\). This is the prolongation-based contact
Hamiltonian picture developed in
\cite{SimoesColomboDeLeonSalgadoSouto2025Annali}, and it is compatible with the
Jacobi description of \(E^*\times\mathbb R\) in
\cite{ASCLMS2024}. The present formulation starts one step earlier: it produces the generated
subset $\mathcal R_L=\operatorname{Im}(\Lambda_L)\subset T(E^*\times\mathbb R)$
directly from the contact Tulczyjew morphism. Therefore it does not require, as
part of the construction, that the relation be the anchor image of a globally
defined or unique dynamical section. This distinction is important for singular
generating objects and for skew algebroids, where we work with the Tulczyjew
morphism and admissibility relation rather than with a full contact Lie
algebroid prolongation.
\end{remark}

%\begin{remark}
%In the singular case the contact Legendre map fails to %be locally invertible, and
%the equations define an implicit relation rather than a vector field. One should
%then expect a contact constraint algorithm, in the spirit of the precontact
%approach to singular Lagrangians \cite{deLeonLainz2019}. We do not develop that
%algorithm here, but the implicit nature of the contact Tulczyjew relation shows
%that singular generating objects are naturally accommodated by the framework.
%\end{remark}


\section{The contact Tulczyjew morphism in line-bundle form}
\label{sec:intrinsic}


\begin{thebibliography}{99}
\bibitem[ASCLMS2024]{ASCLMS2024} A. Anahory Simoes, L. Colombo, M. de Le\'on, M. Salgado and S. Souto, Euler--Lagrange--Herglotz equations on Lie algebroids, \textit{Analysis and Mathematical Physics} \textbf{14} (2024), 3.
\bibitem[GG2008]{GG2008} K. Grabowska and J. Grabowski, Variational calculus with constraints on general algebroids, \textit{Journal of Physics A: Mathematical and Theoretical} \textbf{41} (2008), 175204.
\bibitem[GG2024]{GG2024} K. Grabowska and J. Grabowski, Contact geometric mechanics: the Tulczyjew triple, \textit{Advances in Theoretical and Mathematical Physics} \textbf{28} (2024), 599--654.
\bibitem[Martinez2001]{Martinez2001} E. Mart\'inez, Lagrangian mechanics on Lie algebroids, \textit{Acta Applicandae Mathematicae} \textbf{67} (2001), no. 3, 295--320.
\bibitem[MendellaMarmoTulczyjew1995]{MendellaMarmoTulczyjew1995} G. Mendella, G. Marmo and W. M. Tulczyjew, Integrability of implicit differential equations, \textit{Journal of Physics A: Mathematical and General} \textbf{28} (1995), no. 1, 149--163.
\bibitem[SimoesColomboDeLeonSalgadoSouto2025Annali]{SimoesColomboDeLeonSalgadoSouto2025Annali} A. Anahory Simoes, L. Colombo, M. de Le\'on, M. Salgado and S. Souto, Contact formalism for dissipative mechanical systems on Lie algebroids, \textit{Annali di Matematica Pura ed Applicata} \textbf{204} (2025), 1847--1880.
\bibitem[Tulczyjew1976]{Tulczyjew1976} W. M. Tulczyjew, Les sous-vari\'et\'es lagrangiennes et la dynamique lagrangienne, \textit{Comptes Rendus de l'Acad\'emie des Sciences de Paris, S\'erie A--B} \textbf{283} (1976), no. 8, A675--A678.
\bibitem[deLeonLainz2019]{deLeonLainz2019} M. de Le\'on and M. Lainz Valc\'azar, Singular Lagrangians and precontact Hamiltonian systems, \textit{International Journal of Geometric Methods in Modern Physics} \textbf{16} (2019), no. 10, 1950158.
\end{thebibliography}
\end{document}
